\documentclass[11pt]{article}
\usepackage{amsmath,amssymb,amsthm}
\usepackage[no-math]{fontspec}
\setmainfont{cmunrm.otf}[BoldFont=cmunbx.otf,ItalicFont=cmunti.otf,BoldItalicFont=cmunbi.otf]
\setsansfont{cmunss.otf}[BoldFont=cmunsx.otf,ItalicFont=cmunsi.otf]
\setmonofont{cmuntt.otf}
\usepackage[english,russian]{babel}
\usepackage[a4paper,margin=2.4cm]{geometry}
\usepackage{booktabs,enumitem,xcolor}
\usepackage{forest}
\usepackage[colorlinks=true,linkcolor=blue!50!black,citecolor=blue!50!black,urlcolor=blue!50!black]{hyperref}

\newcommand{\Lie}{\mathrm{Lie}}
\newcommand{\Inv}{\mathrm{Inv}}
\newcommand{\st}{\mathrm{st}}
\newcommand{\cher}{\mathrm{ch}}
\newcommand{\spn}{\mathrm{span}}
\newcommand{\Z}{\mathbb Z}
\newcommand{\Q}{\mathbb Q}
\newcommand{\g}{\gamma}

% метки статуса утверждений
\newcommand{\stat}[1]{\textsf{\small[#1]}}
\newcommand{\proved}{\stat{доказано}}
\newcommand{\comproved}{\stat{доказано с компьютером}}
\newcommand{\computed}{\stat{посчитано}}
\newcommand{\unchecked}{\stat{не перепроверено}}
\newcommand{\open}{\stat{открыто}}

\theoremstyle{plain}
\newtheorem{theorem}{Теорема}
\newtheorem{lemma}[theorem]{Лемма}
\newtheorem{proposition}[theorem]{Предложение}
\newtheorem{corollary}[theorem]{Следствие}
\newtheorem*{conjecture}{Гипотеза}
\theoremstyle{definition}
\newtheorem{definition}[theorem]{Определение}
\newtheorem{example}[theorem]{Пример}
\theoremstyle{remark}
\newtheorem{remark}[theorem]{Замечание}

\forestset{bintree/.style={for tree={math content, inner sep=1pt, l=5.5mm, l sep=0pt, s sep=3.2mm, font=\small},
  where n children=0{}{shape=coordinate}}}

\setlist{itemsep=2pt,topsep=4pt}
\sloppy

\begin{document}

\begin{center}
{\LARGE\bfseries Скобочные тождества алгебр Ли}\\[6pt]
{\large Задача 2.6 из «Днестровской тетради» (Ю.\,А.~Бахтурин):\\ постановка, сведения и состояние на 12 октября 2026 г.}\\[8pt]
{\small Рабочая записка проекта. Скрипты и данные: папка \texttt{projects/monomial-identities} репозитория \texttt{seminar-kb}.}
\end{center}

\bigskip

\section*{Коротко}

\textbf{Задача.} Расставим скобки на буквах $x_1,\dots,x_n$ --- получится лиев одночлен, например $[[x_1,x_2],[x_3,x_4]]$. Каждый такой одночлен задаёт тождество (нильпотентность, метабелевость, разрешимость и так далее). Верно ли, что два таких тождества равносильны только тогда, когда скобки расставлены одинаково (с точностью до перестановки сомножителей в скобках)?

\textbf{Что известно на сегодня.}
\begin{itemize}
\item Общего доказательства нет. \open
\item Гипотеза верна для одночленов степени $n\le 12$ над полем характеристики $0$. \comproved
\item Она верна для всех $n$ внутри семейства полинильпотентных форм (следствие теоремы Парфёнова--Бахтурина). Гребёнка отличается от всех остальных форм. Для форм с двумя «вишнями» $[x,x]$ есть рассуждение через алгебру Витта, оно не перепроверено.
\item По тождеству одночлена восстанавливаются: степень, высота дерева, число вишен и число Стралера (последнее --- по лемме Зайцева, характеристика $0$). \proved
\item В прочитанной литературе решения нет ни для алгебр Ли, ни для аналогичной задачи о группах. Для \emph{произвольных} вербальных идеалов сокращение в коммутаторе неверно (пример Зайцева, 1981), так что доказательство обязано использовать специальный вид идеалов $I_\sigma$.
\item Из гипотезы над $\Q$ следует её аналог для групп; вместе с результатом для $n\le12$ это даёт: разные внешние коммутаторные слова веса не больше $12$ задают разные многообразия групп (набросок, замечание~\ref{rem:groups}).
\end{itemize}

Каждое утверждение в записке помечено: \proved{} --- есть полное доказательство; \comproved{} --- доказательство опирается на конечный перебор, который проверяется программой без приближённых вычислений; \computed{} --- результат счёта, не оформленный как доказательство (например, ранги по модулю простого числа); \unchecked{} --- рассуждение есть, но его никто не проверял; \open{} --- неизвестно.

\tableofcontents

\section{Задача на примерах}

\subsection{Одночлены и их формы}

\emph{Скобочный одночлен} --- это выражение, в котором на буквах $x_1,\dots,x_n$ полностью расставлены лиевы скобки и каждая буква встречается один раз. Скобка антикоммутативна, $[a,b]=-[b,a]$, поэтому порядок двух сомножителей внутри скобки меняет только знак. Если стереть буквы и забыть порядок сомножителей, от одночлена останется его \emph{форма} --- бинарное дерево, у которого дети вершины не упорядочены.

При $n=2$ и $n=3$ форма одна. При $n=4$ их две, при $n=5$ --- три:
\begin{center}
\begin{forest} bintree [ [ [ [x][x] ] [x] ] [x] ] \end{forest}\qquad
\begin{forest} bintree [ [ [x][x] ] [ [x][x] ] ] \end{forest}\qquad\qquad\qquad
\begin{forest} bintree [ [ [ [ [x][x] ] [x] ] [x] ] [x] ] \end{forest}\qquad
\begin{forest} bintree [ [ [x][x] ] [ [ [x][x] ] [x] ] ] \end{forest}\qquad
\begin{forest} bintree [ [x] [ [ [x][x] ] [ [x][x] ] ] ] \end{forest}
\end{center}
Число форм с $n$ листьями при $n=1,\dots,12$ равно
\[1,\ 1,\ 1,\ 2,\ 3,\ 6,\ 11,\ 23,\ 46,\ 98,\ 207,\ 451\]
(числа Веддербёрна--Этерингтона).

Два слова понадобятся постоянно. \emph{Гребёнка} $c_n$ --- это форма левонормированного одночлена
\[ [[\dots[[x_1,x_2],x_3],\dots],x_n]. \]
\emph{Вишня} --- вершина, оба ребёнка которой листья, то есть скобка двух букв $[x,x]$.

\subsection{Тождества}

Одночлен $w$ задаёт тождество $w=0$: алгебра Ли $L$ ему удовлетворяет, если $w$ обращается в нуль при любой подстановке элементов $L$ вместо букв. Примеры:
\begin{itemize}
\item $[x_1,x_2]=0$ --- абелевы алгебры;
\item $[[\dots[x_1,x_2],\dots],x_{c+1}]=0$ --- нильпотентные ступени не выше $c$;
\item $[[x_1,x_2],[x_3,x_4]]=0$ --- метабелевы;
\item полное бинарное дерево глубины $d$ --- разрешимые ступени не выше $d$.
\end{itemize}

Говорят, что тождество $v=0$ \emph{следует} из $w=0$, если оно выполнено во всякой алгебре Ли, в которой выполнено $w=0$. Тождества \emph{равносильны}, если каждое следует из другого.

\begin{example}\label{ex:n4}
При $n=4$ из нильпотентности $[[[x_1,x_2],x_3],x_4]=0$ следует метабелевость: $[[a,b],[c,d]]$ лежит в $[L',L']\subseteq\g_4(L)$. Обратно не следует: двумерная алгебра с базисом $e_1,e_2$ и $[e_1,e_2]=e_2$ метабелева, но не нильпотентна. Значит, два тождества степени $4$ не равносильны.
\end{example}

\subsection{Где трудность}

Кажется очевидным, что разные формулы дают разные тождества. Но тождество Якоби
\[ [[a,b],c]=[[a,c],b]+[a,[b,c]] \]
связывает одночлены \emph{разных} форм, и поэтому из одного скобочного тождества следуют другие, внешне на него не похожие (пример~\ref{ex:n4}). Вопрос в том, может ли случиться, что два одночлена разной формы следуют \emph{друг из друга}. Гипотеза утверждает, что не может.

\section{Точная постановка}

Всюду $k$ --- поле характеристики $0$, все алгебры --- над $k$.

\begin{definition}
\emph{Форма} --- элемент свободного коммутативного неассоциативного группоида (магмы) с одним порождающим $x$. Иначе говоря, формы строятся по правилу: $x$ --- форма; если $\alpha,\beta$ --- формы, то $[\alpha,\beta]$ --- форма, причём $[\alpha,\beta]=[\beta,\alpha]$, и других равенств нет. Число вхождений $x$ в форму $\sigma$ обозначается $|\sigma|$.

Для формы $\sigma$ с $|\sigma|=n$ через $w_\sigma(x_1,\dots,x_n)$ обозначается лиев одночлен, полученный расстановкой букв $x_1,\dots,x_n$ в листья $\sigma$ (каждая буква один раз) и выбором порядка сомножителей в каждой скобке.
\end{definition}

Разные расстановки букв и разные порядки сомножителей дают одночлены, отличающиеся переименованием букв и знаком. Тождество $w_\sigma=0$ от этих выборов не зависит.

\begin{conjecture}[Ю.\,А.~Бахтурин, задача~2.6 «Днестровской тетради»]
Тождества алгебр Ли $w_\sigma=0$ и $w_\tau=0$ равносильны тогда и только тогда, когда $\sigma=\tau$.
\end{conjecture}

В одну сторону утверждение тривиально. Содержательно другое: если $\sigma\ne\tau$, то найдётся алгебра Ли, в которой одно из тождеств выполнено, а другое нет.

\begin{remark}[о тексте тетради]\label{rem:text}
В издании 1993~г.\ \cite{DN} задача записана так: «Доказать, что для полилинейных одночленов $[x_1,\dots,x_n]_\sigma=w_\sigma$ и $[x_1,\dots,x_n]_\tau=w_\tau$ (с расстановками скобок $\sigma$ и $\tau$) тождества алгебр Ли $w_\sigma=0$ и $w_\tau=0$ эквивалентны тогда и только тогда, когда $w_\sigma$ и $w_\tau$ равны как элементы свободного коммутативного (неассоциативного) группоида от переменных $x_1,\dots,x_n$».

Буквально это неверно уже при $n=3$. Одночлены $[[x_1,x_2],x_3]$ и $[x_1,[x_2,x_3]]$ --- разные элементы свободного коммутативного группоида от $x_1,x_2,x_3$ (в первом вместе стоят $x_1,x_2$, во втором $x_2,x_3$), а тождества равносильны: одно получается из другого переименованием букв. Сравнивать нужно после отождествления всех букв, то есть в группоиде с одним порождающим. Так же, через свободный коммутативный группоид с одним порождающим, Дж.~Вильсон \cite{Wil} определяет внешние коммутаторные слова в группах (вопрос о различии многообразий он не ставит). Отметки о решении у задачи в тетради нет.
\end{remark}

\begin{remark}[о поле]
В тетради поле не указано. Леммы и теорема о сведении из разделов \ref{sec:ideals} и \ref{sec:finite} верны над любым полем. Все числа (размерности, таблицы следований) и теорема~\ref{thm:n12} относятся к характеристике $0$. В положительной характеристике размерности уже другие: над полем из двух элементов $\dim M_{[F,F'']}=14$ при $n=5$, а не $15$. \computed{} Верна ли гипотеза в положительной характеристике и над $\Z$, не выяснялось.
\end{remark}

\section{Первое сведение: идеалы свободной алгебры}\label{sec:ideals}

Пусть $F$ --- свободная алгебра Ли со счётным множеством свободных порождающих $x_1,x_2,\dots$. Для формы $\sigma$, $|\sigma|=n$, положим
\[ I_\sigma=\spn\{\,w_\sigma(u_1,\dots,u_n)\ :\ u_1,\dots,u_n\in F\,\}. \]

\begin{lemma}\label{lem:ideal} \proved{}
$I_\sigma$ --- идеал алгебры $F$, выдерживающий все подстановки: если $f(x_1,\dots,x_m)\in I_\sigma$, то $f(v_1,\dots,v_m)\in I_\sigma$ для любых $v_i\in F$. (Такие идеалы называются вербальными.)
\end{lemma}
\begin{proof}
Умножение на элемент --- дифференцирование: $[[u,v],y]=[[u,y],v]+[u,[v,y]]$. Индукцией по построению одночлена
\[ [w_\sigma(u_1,\dots,u_n),y]=\sum_{i=1}^n w_\sigma(u_1,\dots,[u_i,y],\dots,u_n). \]
Справа сумма значений того же одночлена, значит, $I_\sigma$ --- идеал. Подстановка в значение --- снова значение. Полилинейность делает рассуждение верным над любым полем.
\end{proof}

\begin{lemma}\label{lem:comm} \proved{}
$I_x=F$, а если $\sigma=[\alpha,\beta]$, то $I_\sigma=[I_\alpha,I_\beta]$ --- взаимный коммутант идеалов, то есть линейная оболочка всех $[a,b]$, $a\in I_\alpha$, $b\in I_\beta$.
\end{lemma}
\begin{proof}
Буквы в двух частях одночлена $w_{[\alpha,\beta]}$ разные, поэтому его значения --- это в точности элементы $[a,b]$, где $a$ --- значение $w_\alpha$, $b$ --- значение $w_\beta$. Их оболочка совпадает с $[I_\alpha,I_\beta]$ по билинейности скобки.
\end{proof}

Итак, идеал формы вычисляется по дереву: \emph{лист --- это $F$, вершина --- коммутант идеалов её детей}. Гребёнка $c_n$ даёт $\g_n(F)$, $n$-й член нижнего центрального ряда; полное дерево глубины $d$ даёт $F^{(d)}$, $d$-й член производного ряда ($F^{(0)}=F$, $F^{(d+1)}=[F^{(d)},F^{(d)}]$); форма $[[x,x],[[x,x],x]]$ даёт $[F',\g_3(F)]$. Такими именами формы и называются ниже.

\begin{lemma}\label{lem:impl} \proved{}
Тождество $w_\tau=0$ следует из $w_\sigma=0$ тогда и только тогда, когда $I_\tau\subseteq I_\sigma$. В частности, тождества равносильны тогда и только тогда, когда $I_\sigma=I_\tau$.
\end{lemma}
\begin{proof}
Для алгебры $L$ обозначим через $I_\sigma(L)$ оболочку значений $w_\sigma$ в $L$; тождество выполнено в $L$ тогда и только тогда, когда $I_\sigma(L)=0$. Тождество от конечного числа букв достаточно проверять на счётно порождённых алгебрах, а всякая такая алгебра есть образ $F$ при некотором гомоморфизме $\varphi$, причём $I_\sigma(L)=\varphi(I_\sigma)$.

Если $I_\tau\subseteq I_\sigma$ и $\varphi(I_\sigma)=0$, то $\varphi(I_\tau)=0$. Обратно, в алгебре $F/I_\sigma$ тождество $w_\sigma=0$ выполнено; если из него следует $w_\tau=0$, то образ $I_\tau$ в $F/I_\sigma$ нулевой, то есть $I_\tau\subseteq I_\sigma$.
\end{proof}

\begin{corollary}[гипотеза на языке идеалов]\label{cor:groupoid} \proved{}
Следующие утверждения равносильны.
\begin{enumerate}[label=(\roman*)]
\item Гипотеза: из равносильности тождеств $w_\sigma=0$, $w_\tau=0$ следует $\sigma=\tau$.
\item Отображение $\sigma\mapsto I_\sigma$ инъективно.
\item Множество идеалов, получаемых из $F$ повторным применением операции $(A,B)\mapsto[A,B]$, есть свободный коммутативный группоид с порождающим $F$.
\item Однозначность разложения: если $A,B,C,D$ --- идеалы вида $I_\sigma$ и $[A,B]=[C,D]$, то $\{A,B\}=\{C,D\}$.
\end{enumerate}
\end{corollary}
\begin{proof}
(i)$\Leftrightarrow$(ii) --- лемма~\ref{lem:impl}. (ii)$\Leftrightarrow$(iii) --- по определению: формы суть элементы свободного коммутативного группоида, а $\sigma\mapsto I_\sigma$ --- гомоморфизм группоидов по лемме~\ref{lem:comm}. (ii)$\Rightarrow$(iv) очевидно. (iv)$\Rightarrow$(ii): число листьев восстанавливается по идеалу --- это наименьшая степень ненулевого элемента $I_\sigma$; пусть $I_\sigma=I_\tau$, $\sigma=[\alpha,\beta]$, $\tau=[\gamma,\delta]$, тогда $[I_\alpha,I_\beta]=[I_\gamma,I_\delta]$, по (iv), скажем, $I_\alpha=I_\gamma$, $I_\beta=I_\delta$, и по индукции по числу листьев $\alpha=\gamma$, $\beta=\delta$.
\end{proof}

Пункт (iii) и объясняет слова «свободный коммутативный группоид» в формулировке задачи.

\begin{remark}[язык многообразий]
Пусть $\mathbf V_\sigma$ --- многообразие, заданное тождеством $w_\sigma=0$. В книге Бахтурина \cite[п.~4.3.8]{Bah} определён \emph{коммутатор} многообразий $[\mathbf U,\mathbf V]$: это многообразие, вербальный идеал которого равен коммутанту вербальных идеалов $\mathbf U$ и $\mathbf V$. По лемме~\ref{lem:comm} $\mathbf V_{[\alpha,\beta]}=[\mathbf V_\alpha,\mathbf V_\beta]$, а $\mathbf V_x$ состоит из нулевой алгебры. Гипотеза: разные формы задают разные многообразия. М.\,В.~Зайцев называет многообразия $\mathbf V_\sigma$ \emph{коммутаторными} \cite{Z98}.
\end{remark}

\section{Второе сведение: конечномерная задача}\label{sec:finite}

\begin{definition}
$\Lie(n)$ --- полилинейная часть свободной алгебры Ли: линейная оболочка скобочных одночленов от $x_1,\dots,x_n$, в каждом из которых каждая буква встречается ровно один раз. Группа $S_n$ действует на $\Lie(n)$ перестановками букв.
\end{definition}

Классические факты (см., например, \cite{Reu}): $\dim\Lie(n)=(n-1)!$; базис образуют левонормированные скобки $[\dots[[x_1,x_{i_2}],x_{i_3}],\dots,x_{i_n}]$ по всем перестановкам $(i_2,\dots,i_n)$ остальных букв; свободная алгебра Ли вкладывается в свободную ассоциативную ($[u,v]=uv-vu$), и координата элемента $\Lie(n)$ при базисной скобке равна его коэффициенту при слове $x_1x_{i_2}\cdots x_{i_n}$.

\begin{definition}
Для формы $\sigma$ с $n$ листьями $M_\sigma\subseteq\Lie(n)$ --- линейная оболочка одночленов $w_\sigma(x_{\pi(1)},\dots,x_{\pi(n)})$ по всем $\pi\in S_n$. Это $S_n$-подмодуль, порождённый одним элементом.
\end{definition}

\begin{theorem}[о сведении]\label{thm:reduction} \proved{}
Пусть $|\sigma|=|\tau|=n$. Равносильны:
\begin{enumerate}[label=(\arabic*)]
\item тождество $w_\tau=0$ следует из $w_\sigma=0$;
\item $w_\tau\in M_\sigma$;
\item $M_\tau\subseteq M_\sigma$.
\end{enumerate}
Если же $|\sigma|\ne|\tau|$, то тождества $w_\sigma=0$ и $w_\tau=0$ не равносильны.
\end{theorem}
\begin{proof}
(2)$\Leftrightarrow$(3): $M_\sigma$ выдерживает перестановки букв, а $M_\tau$ порождён перестановками одного $w_\tau$.

(2)$\Rightarrow$(1): $w_\tau\in M_\sigma\subseteq I_\sigma$; идеал $I_\sigma$ выдерживает подстановки, поэтому все значения $w_\tau$ лежат в $I_\sigma$, то есть $I_\tau\subseteq I_\sigma$, и остаётся применить лемму~\ref{lem:impl}.

(1)$\Rightarrow$(2): по лемме~\ref{lem:impl} $w_\tau\in I_\sigma$, то есть $w_\tau=\sum_j a_j\,w_\sigma(u^{(j)}_1,\dots,u^{(j)}_n)$, где все $u^{(j)}_i$ можно считать одночленами от букв. Каждое слагаемое однородно по каждой букве. Возьмём в этом равенстве компоненту «буквы $x_1,\dots,x_n$ по одному разу, других букв нет». Слева останется $w_\tau$. Справа останутся слагаемые суммарной степени $n$; в них все $u^{(j)}_i$ --- буквы, то есть перестановка $x_1,\dots,x_n$. Такие слагаемые лежат в $M_\sigma$.

Последнее утверждение: все ненулевые элементы $I_\sigma$ имеют степень не меньше $|\sigma|$, а $w_\tau\ne0$ в $F$; поэтому из $w_\tau\in I_\sigma$ следует $|\tau|\ge|\sigma|$, и из равносильности --- равенство степеней.
\end{proof}

\begin{corollary}[гипотеза на языке подпространств] \proved{}
Гипотеза равносильна тому, что при каждом $n$ подпространства $M_\sigma\subseteq\Lie(n)$, отвечающие разным формам с $n$ листьями, попарно различны.
\end{corollary}

Для каждого отдельного $n$ это конечная задача линейной алгебры. Направление легко перепутать: чем \emph{больше} $M_\sigma$, тем \emph{сильнее} тождество $w_\sigma=0$. У гребёнки $M_{c_n}=\Lie(n)$: нильпотентность влечёт все скобочные тождества той же степени.

\begin{example}
$n=4$: $\dim\Lie(4)=6$; у гребёнки $M=\Lie(4)$, у формы $[[x,x],[x,x]]$ подпространство натянуто на три независимых элемента $[[x_1,x_2],[x_3,x_4]]$, $[[x_1,x_3],[x_2,x_4]]$, $[[x_1,x_4],[x_2,x_3]]$.

$n=5$: $\dim\Lie(5)=24$; размерности трёх подпространств равны $24$, $20$, $15$, и
\[ M_{[F,F'']}\subset M_{[F',\g_3]}\subset M_{\g_5}=\Lie(5). \]
Первое включение отвечает на упражнение 4.10.20 из книги Бахтурина \cite{Bah}: «Которое из тождеств $x(yz)uv\equiv0$, $(xyz)uv\equiv0$ влечёт другое? Почему они неэквивалентны?» Произведения там нужно читать как правонормированные, $x(yz)uv=[x,[[y,z],[u,v]]]$ и $(xyz)uv=[[x,[y,z]],[u,v]]$; при левонормированном чтении обе формы оказались бы гребёнками. Итак, из второго тождества следует первое, а обратно --- нет.
\end{example}

Формы с шестью листьями собраны в приложении~\ref{app:atlas}. Там видно, что размерность формы не различает: у $[F',\g_4]$ и $[F,[F,F'']]$ она равна $75$, а сумма этих подпространств имеет размерность $110$. \computed

\section{Что читается по идеалу: инварианты формы}\label{sec:inv}

\emph{Инвариант} --- характеристика дерева, которую можно вычислить, зная только идеал $I_\sigma$ (или подпространство $M_\sigma$). Если тождества форм $\sigma$ и $\tau$ равносильны, все инварианты у них совпадают. Гипотеза будет доказана, если найти столько инвариантов, что по ним восстанавливается дерево.

\subsection{Якоби для идеалов и выводимость}

\begin{lemma}[о трёх идеалах]\label{lem:three} \proved{}
Для любых идеалов $A,B,C$ алгебры Ли
\[ [[A,B],C]\subseteq[[A,C],B]+[A,[B,C]]. \]
\end{lemma}

Это тождество Якоби, применённое к элементам. Следствия: $[\g_i,\g_j]\subseteq\g_{i+j}$, откуда $I_\sigma\subseteq\g_n(F)$ для любой формы с $n$ листьями (гребёнка сильнее всех); при $B=A$ получается $[[A,A],C]\subseteq[[A,C],A]$.

\begin{definition}[шаг Якоби]
Пусть в форме $\sigma$ есть поддерево вида $[[a,b],c]$. Заменив его на $[[a,c],b]$, получим форму $\tau_1$; заменив на $[a,[b,c]]$ --- форму $\tau_2$. По лемме~\ref{lem:three} $I_\sigma\subseteq I_{\tau_1}+I_{\tau_2}$. Пишем $\sigma\preceq\rho$ («$\sigma$ выводится из $\rho$»), если это получается из правил: $\rho\preceq\rho$; если $\tau_1\preceq\rho$ и $\tau_2\preceq\rho$, то $\sigma\preceq\rho$.
\end{definition}

\begin{proposition}\label{prop:sound} \proved{}
Если $\sigma\preceq\rho$, то $I_\sigma\subseteq I_\rho$: из тождества формы $\rho$ следует тождество формы $\sigma$.
\end{proposition}

Например, $[F'',F]\preceq[F',\g_3]$ за один шаг ($a=b=[x,x]$, $c=x$, обе замены дают одну форму).

\subsection{Высота}

\emph{Высота} $h(\sigma)$ --- наибольшая глубина листа. Пусть $z$ --- лист одночлена $T$ и $u_1,\dots,u_k$ --- поддеревья, висящие на пути от $z$ к корню, снизу вверх; тогда $T=\pm[\dots[[z,u_1],u_2],\dots,u_k]$, и слова одночлена $T$, начинающиеся с буквы $z$, --- это в точности слова из $\pm z\,\widetilde u_1\widetilde u_2\cdots\widetilde u_k$, где $\widetilde u$ --- одночлен $u$ с раскрытыми скобками. Стирая первую букву $x_1$ в словах, начинающихся с $x_1$, получаем изоморфизм $\Lie(n)\to A_{n-1}$ на полилинейную часть свободной ассоциативной алгебры от остальных букв, при котором
\begin{equation}\label{eq:spine}
M_\sigma\ \longmapsto\ \sum_{\ell}\ M_{\alpha_1(\ell)}\cdot M_{\alpha_2(\ell)}\cdots M_{\alpha_{k}(\ell)},
\end{equation}
где сумма берётся по листьям $\ell$ формы $\sigma$, $(\alpha_1(\ell),\dots,\alpha_k(\ell))$ --- формы поддеревьев вдоль пути от $\ell$ к корню, а произведения берутся по всем распределениям букв. \proved

\begin{proposition} \proved{}
Высота $h(\sigma)$ восстанавливается по $M_\sigma$: это наименьшее $k$, при котором образ $M_\sigma$ в \eqref{eq:spine} состоит из линейных комбинаций произведений не более чем $k$ лиевых элементов.
\end{proposition}
\begin{proof}
Свободная ассоциативная алгебра есть универсальная обёртывающая свободной алгебры Ли. По теореме Пуанкаре--Биркгофа--Витта произведение $k$ ненулевых лиевых элементов не лежит в оболочке произведений меньшего числа лиевых элементов: его символ в симметрической алгебре есть произведение $k$ ненулевых множителей, а в кольце многочленов нет делителей нуля. Самый глубокий лист даёт в \eqref{eq:spine} произведения ровно $h(\sigma)$ ненулевых лиевых элементов, остальные листья --- не больше.
\end{proof}

\subsection{Число Стралера}

\begin{definition}
Число Стралера: $\st(x)=0$; $\st([\alpha,\beta])=\st(\alpha)+1$, если $\st(\alpha)=\st(\beta)$, и $\max(\st(\alpha),\st(\beta))$ иначе. У гребёнки оно равно $1$, у полного дерева глубины $d$ равно $d$.
\end{definition}

\begin{proposition}\label{prop:strahler} \stat{доказано; точность --- по лемме~12 из \cite{Z81}, характеристика $0$}
$I_\sigma\subseteq F^{(d)}$ тогда и только тогда, когда $d\le\st(\sigma)$. Таким образом, $\st(\sigma)$ --- номер последнего члена производного ряда, содержащего $I_\sigma$.
\end{proposition}
\begin{proof}
Индукцией по дереву доказываем сразу два утверждения: $I_\sigma\subseteq F^{(\st(\sigma))}$ и $I_\sigma\not\subseteq F^{(\st(\sigma)+1)}$. Для листа $I_x=F$ лежит в $F^{(0)}$ и не лежит в $F'$.

Включение. Пусть $I_\alpha\subseteq F^{(p)}$, $I_\beta\subseteq F^{(q)}$, $p\le q$. Если $p=q$, то $[I_\alpha,I_\beta]\subseteq F^{(p+1)}$. Если $p<q$, то $[I_\alpha,I_\beta]\subseteq I_\beta\subseteq F^{(q)}$, так как $I_\beta$ --- идеал.

Невключение --- лемма~12 статьи Зайцева \cite{Z81}: если вербальные идеалы $U$, $V$ таковы, что $U\subseteq F^{(p)}$, $U\not\subseteq F^{(p+1)}$, $V\subseteq F^{(q)}$, $V\not\subseteq F^{(q+1)}$, то $[U,V]\subseteq F^{(r)}$ и $[U,V]\not\subseteq F^{(r+1)}$, где $r=\max(p,q)$ при $p\ne q$ и $r=p+1$ при $p=q$; это в точности правило вычисления числа Стралера, и шаг индукции --- эта лемма для $U=I_\alpha$, $V=I_\beta$. Она выводится из леммы~3 той же статьи (свободная алгебра Ли счётного ранга над полем характеристики $0$):
\begin{quote}
пусть $U,V,M,N$ --- вербальные идеалы, $[U,V]\subseteq M(N)$, $U\not\subseteq M(N)$, $V\not\subseteq M(N)$; тогда $U+V\subseteq N$.
\end{quote}
Здесь $M(N)$ --- вербальный идеал произведения многообразий; при $M=F'$ это $[N,N]$. (В лемме молчаливо предполагается $M\ne F$: при $M=F$ она неверна.) Пусть $p<q$ и $[U,V]\subseteq F^{(q+1)}=[N,N]$ для $N=F^{(q)}$. Ни $U$, ни $V$ не лежат в $F^{(q+1)}$, поэтому $U\subseteq F^{(q)}\subseteq F^{(p+1)}$ --- противоречие. Пусть $p=q$ и $[U,V]\subseteq F^{(p+2)}=[N,N]$ для $N=F^{(p+1)}$. Ни $U$, ни $V$ не лежат в $F^{(p+2)}$, поэтому $U\subseteq F^{(p+1)}$ --- противоречие.
\end{proof}

Лемма~3 Зайцева доказывается через сплетения алгебр Ли А.\,Л.~Шмелькина \cite{Shm}; обе формулировки сверены по тексту статьи \cite{Z81}. Основное поле там характеристики~$0$.

\subsection{Число вишен}

\begin{proposition}\label{prop:cherries} \proved{}
Пусть $\cher(\sigma)$ --- число вишен формы $\sigma$ с $n\ge2$ листьями. Тогда $n-\cher(\sigma)$ равно наибольшему $m$, для которого в $I_\sigma$ есть ненулевой полиоднородный элемент степени $n$, имеющий степень $m$ по одной из букв. В частности, $\cher(\sigma)$ --- инвариант.
\end{proposition}
\begin{proof}
Компонента степени $n$ идеала $I_\sigma$ есть линейная оболочка значений $w_\sigma$ на буквах (возможно, повторяющихся), и каждое такое значение полиоднородно. Поэтому полиоднородный элемент степени $n$ из $I_\sigma$, имеющий степень $m$ по букве $y$, есть линейная комбинация значений, в которых $y$ стоит ровно в $m$ листьях, а в остальных --- любые другие буквы. Если $m>n-\cher(\sigma)$, то листьев без $y$ меньше, чем вишен; значит, в какой-то вишне оба листа --- $y$, и каждое такое значение равно нулю.

Пусть теперь $y$ стоит во всех листьях, кроме одного листа в каждой вишне, а в этих оставшихся листьях стоят попарно разные буквы, отличные от $y$; листьев с буквой $y$ тогда $n-\cher(\sigma)$. Тогда всякое поддерево с двумя и более листьями содержит вишню, а значит, букву, не встречающуюся вне этого поддерева. Покажем индукцией, что значение каждого поддерева ненулевое. В свободной алгебре Ли $[a,b]=0$ только для пропорциональных $a,b$ (подалгебра, порождённая $a$ и $b$, свободна по теореме Ширшова--Витта, а абелева свободная алгебра Ли одномерна). Значения двух соседних поддеревьев не пропорциональны: либо хотя бы одно из них содержит букву, которой нет в другом, либо оба --- листья с буквой $y$, а это исключено выбором.
\end{proof}

\begin{remark}
В характеристике $0$ то же следует из леммы~1 статьи \cite{Z98}: для вербальных идеалов $U,V,I$ с $V\subseteq I$ выполнено $d_I([U,V])=d_I(U)+d_I(V)$, где $d_I(S)$ --- наибольшее $m$, при котором $S$ лежит в $m$-й степени идеала $I$ (и $0$, если $S\not\subseteq I$). При $I=F'$ получается $d_{F'}(I_\sigma)=\cher(\sigma)$. Другие идеалы $I$ дают целое семейство инвариантов того же рода.
\end{remark}

\subsection{Пример, на котором все инварианты молчат}

При $n=6$ формы $P=[F',\g_4]=[[x,x],[x,[x,[x,x]]]]$ и $Q=[F,[F,F'']]=[x,[x,[[x,x],[x,x]]]]$ имеют одинаковые размерность ($75$), высоту ($4$), число Стралера ($2$) и число вишен ($2$). Подпространства при этом разные. Наборы поддеревьев вдоль пути от самого глубокого листа у них одинаковы, $\{x,x,x,[x,x]\}$, но идут в разном порядке; символ из теоремы Пуанкаре--Биркгофа--Витта коммутативен и порядка не видит. Нужны более тонкие инварианты.

\section{Тонкие алгебры и доказательство для $n\le12$}\label{sec:thin}

\begin{definition}
Алгебра Ли $T$ называется \emph{тонкой}, если у неё есть базис $\{e_g\}$, занумерованный элементами абелевой группы $G$, в котором
\[ [e_g,e_h]=c(g,h)\,e_{g+h},\qquad c(g,h)\in k. \]
\end{definition}

Три семейства таких алгебр:
\begin{itemize}
\item \emph{алгебра Витта}: $G=\Z$, $c(i,j)=j-i$;
\item \emph{торические}: $G=\Z^r$, $c(u,v)=\omega(u,v)$ для кососимметричной билинейной целочисленной формы $\omega$ (тождество Якоби проверяется прямым вычислением);
\item \emph{скобка Пуассона}: мономы Лорана от двух переменных со скобкой $\{f,g\}=f_xg_y-f_yg_x$; в базисе $e_p=x^{p_1+1}y^{p_2+1}$, $p\in\Z^2$, получается $c(p,q)=\det\bigl(p+(1,1),\,q+(1,1)\bigr)$.
\end{itemize}

Зафиксируем тонкую алгебру и \emph{цвета} $g_1,\dots,g_r\in G$. Раскраска листьев формы $\sigma$ в $r$ цветов означает подстановку базисного элемента $e_{g_a}$ в каждый лист цвета $a$. Значение одночлена на базисных элементах вычисляется по дереву:
\begin{equation}\label{eq:value}
w_\sigma(e_{\ldots})=\Bigl(\prod_{v} c\bigl(g(L_v),g(R_v)\bigr)\Bigr)\,e_{g(\sigma)},
\end{equation}
где произведение берётся по внутренним вершинам $v$, $L_v$ и $R_v$ --- поддеревья двух детей $v$, а $g(S)$ --- сумма цветов листьев поддерева $S$. Значение ненулевое тогда и только тогда, когда ни в одной вершине множитель не равен нулю.

\emph{Состав} раскраски --- набор $m=(m_1,\dots,m_r)$, где $m_a$ --- число листьев цвета $a$.

\begin{definition}
$\Inv(\sigma)$ --- множество составов $m$, для которых существует раскраска формы $\sigma$ с этим составом и ненулевым значением \eqref{eq:value}. (Это множество зависит от алгебры и выбора цветов.)
\end{definition}

\begin{lemma}[монотонность]\label{lem:mono} \proved{}
Если $M_\sigma\subseteq M_\tau$, то $\Inv(\sigma)\subseteq\Inv(\tau)$ для любой тонкой алгебры и любых цветов. В частности, если $\Inv(\sigma)\ne\Inv(\tau)$ хотя бы для одной тонкой алгебры, то тождества $w_\sigma=0$ и $w_\tau=0$ не равносильны.
\end{lemma}
\begin{proof}
По условию в свободной алгебре Ли
\[ w_\sigma(x_1,\dots,x_n)=\sum_{\pi\in S_n} a_\pi\,w_\tau(x_{\pi(1)},\dots,x_{\pi(n)}). \]
Пусть $m\in\Inv(\sigma)$; подставим вместо букв базисные элементы согласно раскраске с составом $m$ и ненулевым значением. Слева ненулевой элемент, значит, ненулевое и какое-то слагаемое справа. Оно есть значение $w_\tau$ на тех же базисных элементах в другом порядке, то есть на раскраске формы $\tau$ с тем же составом. Поэтому $m\in\Inv(\tau)$.
\end{proof}

\begin{example}\label{ex:witt5}
Алгебра Витта, цвета --- целые числа (веса листьев); значение ненулевое, когда ни в одной вершине суммы весов двух поддеревьев не совпадают. Возьмём набор весов $\{-1,-1,1,1,1\}$ и две формы с пятью листьями.

Форма $[F',\g_3]=[[a,b],[c,[d,e]]]$: положим $a=-1$, $b=1$, $d=-1$, $e=1$, $c=1$. В вишнях веса разные; в вершине $[c,[d,e]]$ суммы $1$ и $0$; в корне суммы $0$ и $1$. Значение ненулевое.

Форма $[F,F'']=[a,[[b,c],[d,e]]]$: в каждой из двух вишен веса должны быть разными, то есть каждая вишня --- это $\{-1,1\}$; но тогда суммы двух вишен равны, и вершина $[[b,c],[d,e]]$ даёт нуль. Подходящей раскраски нет.

По лемме~\ref{lem:mono} $M_{[F',\g_3]}\not\subseteq M_{[F,F'']}$: из тождества $[x,[[x,x],[x,x]]]=0$ не следует тождество $[[x,x],[x,[x,x]]]=0$. Это ответ на вторую половину упражнения 4.10.20, полученный без вычисления размерностей.
\end{example}

\begin{theorem}\label{thm:n12} \comproved{}
Над полем характеристики $0$ гипотеза верна для всех форм с $n\le12$ листьями: если $\sigma\ne\tau$, $|\sigma|=|\tau|\le12$, то тождества $w_\sigma=0$ и $w_\tau=0$ не равносильны.
\end{theorem}
\begin{proof}
При $|\sigma|\ne|\tau|$ это теорема~\ref{thm:reduction}. Для каждого $n\le12$ предъявлен конечный список тонких алгебр с наборами цветов (при $n=11$ и $n=12$ это одни и те же девять пар из приложения~\ref{app:cert}) со свойством: для любых двух разных форм с $n$ листьями хотя бы у одной пары списка множества $\Inv$ различны. Остаётся применить лемму~\ref{lem:mono} и теорему~\ref{thm:reduction}. Множества $\Inv$ вычисляются перебором раскрасок по формуле \eqref{eq:value} в целых числах; линейная алгебра и приближённые вычисления не используются. Проверка: \texttt{checks/thin\_search.py verify}.
\end{proof}

При $n=12$ форм $451$; проверка занимает несколько минут. Алгебры найдены случайным поиском среди торических алгебр и скобок Пуассона с небольшими целыми параметрами.

\begin{remark}
Независимо от этого, для $n\le8$ подпространства $M_\sigma$ сравнивались напрямую, через ранги: все они попарно различны. \computed{} Размерности $\dim M_\sigma$ при $n\le6$ посчитаны точно, при $n=7,8$ --- по модулю простых $2^{61}-1$ и $2^{31}-1$; включения между подпространствами проверялись рангами по модулю $2^{61}-1$. Этот счёт даёт больше, чем теорема~\ref{thm:n12}: полную таблицу следований (приложение~\ref{app:atlas}). Сертификат из тонких алгебр видит не все неследования: при $n=8$ --- $365$ из $416$.
\end{remark}

\section{Одна алгебра Витта: что она умеет и где её не хватает}\label{sec:witt}

Естественно спросить, хватит ли одной явной тонкой алгебры для всех $n$. Для алгебры Витта $\Inv_W(\sigma)$ --- это множество наборов из $n$ целых чисел, которые можно расставить в листья $\sigma$ так, чтобы ни в одной вершине суммы по двум поддеревьям не совпали.

\begin{itemize}
\item \textbf{Правило склейки.} $\Inv_W([\alpha,\beta])$ состоит из объединений $m_1\sqcup m_2$, где $m_1\in\Inv_W(\alpha)$, $m_2\in\Inv_W(\beta)$ и суммы чисел в $m_1$ и $m_2$ различны. \proved{} Тем самым $\sigma\mapsto\Inv_W(\sigma)$ --- гомоморфизм группоидов, и вопрос снова в его инъективности.
\item \textbf{Формы с одной и двумя вишнями.} Одна вишня --- это гребёнка; она отличается от всех остальных форм уже по предложению~\ref{prop:cherries}. \proved{} Формы с двумя вишнями, $[x,[x,\dots[c_a,c_b]\dots]]$, распознаются алгеброй Витта среди всех форм при любом $n$. \unchecked
\item \textbf{При $n\le10$} одной алгебры Витта достаточно: восемь наборов цветов различают все формы. \comproved{} (\texttt{checks/witt/witt\_check.py}.)
\item \textbf{При $n=11$} формы
\[ A=[x,[c_2,[c_2,[c_2,c_4]]]],\qquad B=[x,[c_2,[c_2,[x,[c_2,c_3]]]]] \]
имеют одинаковые множества допустимых наборов для всех весов от $-6$ до $6$: по $1\,345\,542$ набора с каждой стороны. \computed{} Доказательства, что $\Inv_W(A)=\Inv_W(B)$ при любых весах, нет. \open{} Пару различают торическая алгебра и скобка Пуассона из приложения~\ref{app:cert}.
\end{itemize}

Вывод: одна алгебра Витта полным инвариантом, судя по всему, не является. Если это подтвердится, естественные кандидаты для общего доказательства --- тонкие алгебры с «двумерными» весами: кососимметричная форма ранга не меньше $2$ или скобка Пуассона. Другие тонкие алгебры с группой $\Z$ не рассматривались. Одна алгебра Пуассона, без торических, различает все формы при $n\le9$ (при $n=9$ --- восемью наборами точек; дальше не считалось). \computed

\section{Порядок следований и шаги Якоби}

Счёт рангов при $n\le8$ даёт полную таблицу следований между формами одной степени. Число следований между разными формами при $n=4,\dots,8$ равно $1,3,9,28,90$.

\begin{itemize}
\item \textbf{Наблюдение J1.} При $n\le8$ каждое включение $M_\sigma\subseteq M_\rho$ выводится шагами Якоби: $\sigma\preceq\rho$. \computed
\item \textbf{Наблюдение J2.} При $n\le17$ (при $n=17$ это $24\,631$ форма) нет двух разных форм, каждая из которых выводится из другой. \computed{} Это чистая комбинаторика деревьев.
\end{itemize}

Если бы J1 и J2 были верны при всех $n$, гипотеза следовала бы немедленно: из $M_\sigma=M_\tau$ по J1 каждая форма выводится из другой, и по J2 они совпадают. Оба утверждения для общего $n$ открыты. \open{} Вместе они сильнее гипотезы: J1 описывает все следования, а не только отсутствие равносильностей. Само J2 --- необходимое следствие гипотезы: две разные формы, выводимые друг из друга, были бы контрпримером (предложение~\ref{prop:sound}).

Что известно против простых планов доказательства:
\begin{itemize}
\item Для \emph{систем} скобочных тождеств шагов Якоби не хватает. При $n=7$ есть набор из шести форм, замкнутый относительно шагов Якоби, из тождеств которого следует тождество формы $[x,[x,[x,[[x,x],[x,x]]]]]$, хотя ни из одной формы набора по отдельности оно не следует. \computed{} Поэтому J1, если оно верно, нельзя доказать рассуждением о суммах и пересечениях подпространств.
\item Пересечения $M_\sigma$ со слоями фильтрации Пуанкаре--Биркгофа--Витта больше, чем даёт подсчёт по глубине листьев, уже при $n=5$. \computed
\item Инварианты «общего положения» (графы-деревья, раскраски с алгебраически независимыми весами, общая симплектическая форма) различают все формы при $n\le6$ и перестают при $n=7$. \computed{} Сертификат из раздела~\ref{sec:thin} использует специальные целые параметры.
\end{itemize}

\section{Что есть в литературе}

Прочитано по тексту источника.

\textbf{Сама задача.} В русском издании тетради \cite{DN} задача 2.6 стоит без отметки о решении. В книге Бахтурина \cite{Bah} коммутатор многообразий определён в п.~4.3.8; теорем о нём в \S\,4.3 нет, сказано, что коммутатор двух различных неразложимых многообразий «как правило» неразложим относительно умножения, со ссылкой на \cite{Z81} и на упражнение 4.10.14 (в характеристике $0$ многообразие $[\mathbf V,\mathbf E]$ неразложимо; на языке форм это формы вида $[x,\alpha]$). В упражнениях есть ещё 4.10.15\,в) $[\mathbf U_1,\mathbf U_2]\mathbf V=[\mathbf U_1\mathbf V,\mathbf U_2\mathbf V]$, 4.10.17 и 4.10.20 --- случай $n=5$ нашей задачи. В книге прочитаны \S\,4.3, начало \S\,4.4, упражнения и комментарий к главе~4; с.~158--163, 375 и остальные главы не прочитаны. В обзоре \cite{BO} (с.~212) коммутирование многообразий упомянуто двумя фразами, без утверждений о единственности. Зайцев в \cite{Z98} вводит термин «коммутаторное многообразие» и, говоря о базисном ранге, замечает, что класс многообразий, задаваемых одним полилинейным одночленным тождеством, «не был исследован»; о различии таких многообразий там ничего не утверждается.

\textbf{Образец: умножение многообразий.} Многообразия алгебр Ли над бесконечным полем образуют относительно умножения свободную полугруппу с нулём и единицей, её свободные порождающие --- неразложимые многообразия (\cite[теорема 4.3.7]{Bah}; по обзору \cite[с.~210]{BO}, В.\,А.~Парфёнов \cite{Par} доказал это для характеристики $0$, Ю.\,А.~Бахтурин \cite{Bah73} --- для любого бесконечного поля). Ключевой шаг --- теорема о монотонности \cite[4.3.3]{Bah}: если $\mathbf V(K)\subseteq\mathbf V(H)$ для идеалов $H,K$ неабелевой свободной алгебры Ли и многообразия $\mathbf V$, отличного от многообразия всех алгебр Ли, то $K\subseteq H$. Задача 2.6 --- тот же вопрос для другой операции, коммутатора.

\begin{corollary}\label{cor:polynil} \stat{вывод из теоремы 4.3.7; в источниках так не сформулировано}
Назовём форму полинильпотентной, если она получается из гребёнок $c_a$, $a\ge2$, повторной подстановкой гребёнки во все листья гребёнки. Две разные полинильпотентные формы задают неравносильные тождества.
\end{corollary}
\begin{proof}
Идеал такой формы есть $\g_{a_l}(\dots\g_{a_1}(F)\dots)$ --- вербальный идеал произведения нильпотентных многообразий $\mathbf N_{a_l-1}\cdots\mathbf N_{a_1-1}$, все $a_i\ge2$. Нильпотентные многообразия неразложимы \cite[п.~4.3.8]{Bah}, поэтому по теореме 4.3.7 разным последовательностям $(a_1,\dots,a_l)$ отвечают разные многообразия; в частности, разные полинильпотентные формы задаются разными последовательностями.
\end{proof}

\textbf{Коммутатор: чего ждать нельзя.} Для коммутатора аналог теоремы о монотонности неверен. У Зайцева \cite[\S\,2]{Z81} над любым полем построен вербальный идеал $U$, не содержащий $F^{(4)}$, для которого
\[ [U,F']=[U+F^{(4)},F']. \]
Тем самым для произвольных вербальных идеалов из $[A,C]=[B,C]$ и $A\subseteq B$ не следует $A=B$, и группоид всех многообразий относительно коммутатора не свободен. Гипотеза утверждает сокращение только для идеалов вида $I_\sigma$, и её доказательство обязано использовать их специальный вид.

Тот же пример показывает, что и внутри семейства $I_\sigma$ нельзя сокращать \emph{включения}. Пусть $\rho=[[[[c_2,c_2],x],c_2],x]$ (восемь листьев). По лемме~\ref{lem:three}
\[ [F''',F]\subseteq[[F'',F],F'']\subseteq[[[F'',F],F'],F']\subseteq[[[[F'',F],F'],F],F]=[I_\rho,F], \]
хотя $F'''\not\subseteq I_\rho$ \cite[\S\,2, предложение~1]{Z81}; с нашим счётом при $n=8$ это согласуется ($M_{F'''}\not\subseteq M_\rho$, $\dim M_{F'''}=315$). \computed{} Итак, из $[A,F]\subseteq[B,F]$ не следует $A\subseteq B$ даже для $A=I_\alpha$, $B=I_\beta$.

\textbf{Коммутатор: что есть.} У Зайцева \cite{Z81}, кроме лемм 3 и 12, есть лемма~8: если $U\subseteq U'$, $V\subseteq V'$ --- ненулевые вербальные идеалы и $[U,V]=[U',V']$, то $U+V=U'+V'$. Пример выше показывает, что это сокращение «с точностью до суммы» усилить нельзя. Лемма~1 из \cite{Z98} приведена в замечании после предложения~\ref{prop:cherries}.

\textbf{Группы.} Аналог задачи --- вопрос, задают ли разные \emph{внешние коммутаторные слова} разные многообразия групп. Вильсон \cite{Wil} рассматривает ровно свободный коммутативный группоид с одним порождающим и его гомоморфизм в группоид «подгрупповых функций»; вопрос об инъективности он не ставит. Н.\,Р.~Брумберг \cite{Bru} в 1969~г.\ поставил задачу~3.11: пусть $U\subseteq W$, $V\subseteq Z$ --- вербальные подгруппы и $[U,V]=[W,Z]$; доказать или опровергнуть, что $U=W$, $V=Z$. Для цепочки $U\subseteq V\subseteq W$ он доказал, что из $[U,V]=[U,W]$ следует $V=W$ (теорема~2.1). В общем виде ответ отрицательный: пример из \cite[\S\,3]{Z81} --- перенос на группы примера, приведённого выше. В книге Х.~Нейман \cite{Neu} коммутант многообразий определён, утверждения о единственности нет. В работе \cite{FAM} внешние коммутаторные слова изображаются бинарными деревьями; о различии многообразий там ничего.

\begin{remark}\label{rem:groups} \stat{наше наблюдение, набросок}
Из гипотезы над $\Q$ следует её аналог для групп. Пусть $\Phi$ --- свободная группа и $w_\sigma(\Phi)$ --- вербальная подгруппа внешнего коммутаторного слова формы $\sigma$ с $n$ листьями. Присоединённое кольцо Ли группы $\Phi$ свободно (теорема Магнуса--Витта), а групповой коммутатор индуцирует на факторах нижнего центрального ряда его скобку. Поэтому $w_\sigma(\Phi)\subseteq\g_n(\Phi)$, и образ $w_\sigma(\Phi)$ в $\g_n(\Phi)/\g_{n+1}(\Phi)$ есть $\Z$-оболочка значений лиева одночлена $w_\sigma$ на элементах степени $1$; он ненулевой, так что число листьев восстанавливается по подгруппе. Если $w_\sigma(\Phi)=w_\tau(\Phi)$, то эти образы совпадают; после тензорного умножения на $\Q$ это компоненты степени $n$ идеалов $I_\sigma$ и $I_\tau$, а их полилинейные части --- $M_\sigma$ и $M_\tau$.

Вместе с теоремой~\ref{thm:n12} отсюда получается: разные внешние коммутаторные слова веса не больше $12$ задают разные многообразия групп.
\end{remark}

\textbf{Не проверено.} Базы MathSciNet и zbMATH; примечания о решённых задачах в английском издании тетради (2006); обзор Ю.\,А.~Бахтурина «Some thoughts on the current state of the theory of identical relations in Lie algebras» (2021); статья А.\,А.~Урмана «Группоид многообразий некоторых алгебр» (Алгебра и логика, 1969); статья Шмелькина \cite{Shm} известна по ссылкам; работа \cite{FAM} прочитана в виде препринта arXiv:0911.3048.

\section{Что осталось и что делать дальше}

\begin{enumerate}
\item \textbf{Общее доказательство.} Главный кандидат --- тонкие алгебры с двумерными весами: найти явное семейство (торическое или скобку Пуассона с параметрами) и доказать, что по $\Inv([\alpha,\beta])$ восстанавливается пара $\{\alpha,\beta\}$. Первый шаг --- описать для кососимметричной формы, какие вершины дерева вырождаются, по образцу анализа для алгебры Витта.
\item \textbf{Однозначность разложения напрямую.} Образец --- умножение многообразий, но прямого аналога теоремы о монотонности для коммутатора нет: сокращение неверно для произвольных вербальных идеалов, а включения нельзя сокращать даже для идеалов вида $I_\sigma$ (раздел о литературе). Остаются леммы 3 и 8 Зайцева и лемма~1 из \cite{Z98}. Конкретный вопрос: хватает ли инвариантов, получаемых из них при разных идеалах $M$, $N$, $I$, чтобы различить пары форм, на которых молчат высота, число Стралера и число вишен.
\item \textbf{Проверить} рассуждения про алгебру Витта (формы с двумя вишнями) и либо доказать, либо опровергнуть равенство $\Inv_W(A)=\Inv_W(B)$ для пары при $n=11$.
\item \textbf{J1 при $n=9$--$11$} через неприводимые представления $S_n$; комбинаторное описание порядка $\preceq$ и доказательство J2.
\item \textbf{Дочитать}: с.~158--163 и 375 книги \cite{Bah}, статью Урмана, обзор Бахтурина 2021~г.; проверить MathSciNet и zbMATH.
\item \textbf{Положительная характеристика и $\Z$.}
\end{enumerate}

\appendix

\section{Формы с шестью листьями}\label{app:atlas}

\begin{center}
\begin{tabular}{@{}llrl@{}}
\toprule
форма & идеал & $\dim M_\sigma$ & из тождества каких форм следует \\
\midrule
$[x,[x,[x,[x,[x,x]]]]]$ & $\g_6$ & $120$ & --- \\
$[x,[[x,x],[x,[x,x]]]]$ & $[F,[F',\g_3]]$ & $105$ & $\g_6$ \\
$[[x,x],[x,[x,[x,x]]]]$ & $[F',\g_4]$ & $75$ & $\g_6$ \\
$[x,[x,[[x,x],[x,x]]]]$ & $[F,[F,F'']]$ & $75$ & $\g_6$, $[F,[F',\g_3]]$ \\
$[[x,[x,x]],[x,[x,x]]]$ & $[\g_3,\g_3]$ & $40$ & $\g_6$ \\
$[[x,x],[[x,x],[x,x]]]$ & $[F',F'']$ & $30$ & все, кроме $[\g_3,\g_3]$ \\
\bottomrule
\end{tabular}
\end{center}
Здесь $\dim\Lie(6)=120$, а следований между разными формами всего девять. Размерности посчитаны точно, включения --- рангами по модулю простого $2^{61}-1$; данные лежат в файле \texttt{checks/atlas.json}. \computed

\section{Сертификат для $n=11$ и $n=12$}\label{app:cert}

Девять пар «алгебра и набор цветов» --- три торические алгебры и шесть наборов точек в алгебре Пуассона, --- множества $\Inv$ которых различают все $207$ форм с $11$ листьями и все $451$ форму с $12$ листьями.

\emph{Торические}: $G=\Z^r$, цвета --- базисные векторы, $c(u,v)=u^{\mathsf T}\Omega v$.
\[
\Omega_1=\begin{pmatrix}0&-1&0\\1&0&-1\\0&1&0\end{pmatrix},\quad
\Omega_2=\begin{pmatrix}0&1&-2&0\\-1&0&2&2\\2&-2&0&0\\0&-2&0&0\end{pmatrix},\quad
\Omega_3=\begin{pmatrix}0&-2&0&1&-1\\2&0&-1&0&-2\\0&1&0&0&0\\-1&0&0&0&1\\1&2&0&-1&0\end{pmatrix}.
\]

\emph{Алгебра Пуассона}: цвету $a$ отвечает точка $p_a\in\Z^2$, то есть моном $x^{p_{a,1}+1}y^{p_{a,2}+1}$. Шесть наборов точек:
\begin{align*}
&(1,1),\ (-1,-2),\ (1,-2),\ (1,1); &&(-2,1),\ (0,-1),\ (2,-2),\ (0,-2);\\
&(1,-1),\ (1,1),\ (-1,0),\ (2,0); &&(-1,1),\ (-2,1),\ (0,2),\ (2,-1);\\
&(-1,2),\ (1,-1),\ (1,0),\ (-1,1); &&(1,2),\ (0,2),\ (1,2),\ (1,-2),\ (1,0).
\end{align*}
Одинаковые точки в одном наборе --- разные цвета с одним и тем же базисным элементом; на множество $\Inv$ это влияет через состав. Источник: \texttt{checks/thin\_certificates.json}. Сертификаты для $n\le10$ содержат от одной до восьми таких пар.

\section{Где что лежит и как проверить}

\begin{itemize}
\item \texttt{project.md} --- страница проекта; \texttt{START-HERE.md} --- состояние работы и следующие шаги; \texttt{lecture.html} --- курс с вопросами и подсказками.
\item \texttt{checks/thin\_search.py verify} --- проверка теоремы~\ref{thm:n12}; нужен только Python.
\item \texttt{checks/witt/} --- счёт по алгебре Витта: \texttt{witt\_check.py 10}, \texttt{separate\_pair.py}.
\item \texttt{checks/bahturin26.py}, \texttt{atlas.py} --- ранги и таблица следований при $n\le8$; \texttt{jacobi\_order.py}, \texttt{jacobi\_antisymmetry.py} --- наблюдения J1 и J2; \texttt{lattice.py}, \texttt{filtration.py} --- отрицательные результаты. Всем шести нужен пакет \texttt{python-flint}.
\end{itemize}

\begin{thebibliography}{99}
\bibitem{DN} Днестровская тетрадь: нерешённые проблемы теории колец и модулей. 4-е изд. Новосибирск, 1993.
\bibitem{Bah} Бахтурин Ю.\,А. Тождества в алгебрах Ли. М.: Наука, 1985.
\bibitem{BO} Бахтурин Ю.\,А., Ольшанский А.\,Ю. Тождества // Итоги науки и техники. Современные проблемы математики. Фундаментальные направления. Т.~18. М.: ВИНИТИ, 1988. С.~117--240.
\bibitem{Z81} Зайцев М.\,В. О разложимости в произведение коммутаторов многообразий алгебр Ли и групп // Матем. сб. 1981. Т.~116(158), №\,3. С.~315--330.
\bibitem{Z98} Зайцев М.\,В. Суперранг многообразий алгебр Ли // Алгебра и логика. 1998. Т.~37, №\,4. С.~394--412.
\bibitem{Par} Парфёнов В.\,А. О многообразиях алгебр Ли // Алгебра и логика. 1967. Т.~6, №\,4. С.~61--73.
\bibitem{Bah73} Бахтурин Ю.\,А. О тождествах в алгебрах Ли. I, II // Вестник МГУ. Сер. матем., мех. 1973. №\,1. С.~12--18; №\,2. С.~26--33.
\bibitem{Shm} Шмелькин А.\,Л. Сплетения алгебр Ли и их применение в теории групп // Труды ММО. 1973. Т.~29. С.~247--260.
\bibitem{Bru} Брумберг Н.\,Р. О взаимном коммутанте двух многообразий групп // Матем. сб. 1969. Т.~79(121), №\,1. С.~37--58.
\bibitem{Neu} Нейман Х. Многообразия групп. М.: Мир, 1969.
\bibitem{Wil} Wilson J. On outer-commutator words // Canad. J. Math. 1974. Vol.~26. P.~608--620.
\bibitem{FAM} Fern\'andez-Alcober G.\,A., Morigi M. Outer commutator words are uniformly concise // J. London Math. Soc. 2010. Vol.~82. P.~581--595.
\bibitem{Reu} Reutenauer C. Free Lie algebras. Oxford: Clarendon Press, 1993.
\end{thebibliography}

\medskip
{\small Выходные данные \cite{BO}, \cite{Par}, \cite{Bah73}, \cite{Shm}, \cite{FAM}, \cite{Reu} приведены по спискам литературы прочитанных источников и по памяти; перед цитированием в статье их нужно сверить. Известные расхождения: год тома \cite{BO} --- 1987 на реферативной карточке и 1988 в базах; страницы второй части \cite{Bah73} --- 26--33 в списке литературы книги \cite{Bah} и 30--36 в списке статьи \cite{Z81}.}

\end{document}
